\section{The wings belong to the field, and the ratio measured}
\label{sec:ssrfield}

The last section left a tie to break: two closed forms, one frozen
field, and only a computation can adjudicate.  This section runs the
computation --- and finds that the question was subtler than the tie.
The experiment needs fields whose shape we control exactly, so it runs
in the synthetic world, with every construction stated (grid and
solver details in appendix~\ref{app:ssrprotocol}).

Take \emph{two} time-homogeneous fields, built to agree at the money:
both have volatility $20\%$ at the forward, and both have log-slope
$b=\MacSsrHalfSlopeLoc$ there.  They differ only in what ``a straight
skew'' means:
\begin{align*}
\sqrt{v}(y)&=0.20+b\,\log y
&&\text{(straight in \emph{log}-strike)},\\
\sqrt{v}(y)&=0.20+b\,(y-1)
&&\text{(straight in \emph{dollar} strike)}.
\end{align*}
\Cref{fig:ssrreprice}(a) draws both (dashed) with the smiles each
produces today at $\tau=\MacSsrRepTcmp$ (solid), priced through the
forward equation as before.  Walk the panel from the money outward.
At $k=0$ the four curves agree by construction.  The dashed fields
separate quickly --- straight in one chart is curved in the other ---
but the two \emph{solid} curves stay nearly on top of each other
across the whole span: the smiles differ only through second-order
averaging effects, a few tens of basis points at the edges.  Squint,
and today's market looks the same under either field.  Chapter~4
proved this is no accident: many sheets reprice the same quotes.  The
two fields here are a minimal, fully explicit instance of that
ambiguity.

\begin{figure}[htbp]
\centering
\paperfig{\textwidth}{fig_ssr_reprice.pdf}
\caption{The reprice adjudicates (synthetic fields, forward-equation
pricing throughout).  (a)~Two fields (dashed) sharing ATM value and
ATM slope, one straight in log-strike, one in dollar strike, and
their nearly indistinguishable smiles today at
$\tau=\MacSsrRepTcmp$ (solid).  (b)~Each field frozen, the forward
moved $-\MacSsrRepRespMovePct\%$, both surfaces repriced: the smile
changes agree at the money (\MacSsrRepAtmLogBp\ and
\MacSsrRepAtmDlrBp\ vol bp --- the half-rule) and part in the wings,
by \MacSsrRepSepPutBp\ vol bp at $k=-0.30$ and \MacSsrRepSepCallBp\
at $k=+0.30$: flat response against tilted response.  (c)~The
realized ratio $\widehat{\mathcal{R}}(\tau)$ of the reprice: the
straight field pins $2.00$ at every maturity (spread
\MacSsrRepFlatSpread); bending the field lifts the ratio to the
computable level $2+2cH/b$ (dotted guides): \MacSsrRepBentMean\
measured against \MacSsrRepBentPred\ predicted at
$H=-\MacSsrRepRatioMovePct\%$, \MacSsrRepBentBigMean\ against
\MacSsrRepBentBigPred\ at $H=-4\%$.}
\label{fig:ssrreprice}
\end{figure}

Now freeze each field and run the same scenario: forward down
$\MacSsrRepRespMovePct\%$, reprice, and record how the smile
\emph{changed} --- $\sigma_{\mathrm{new}}(k)-\sigma_{\mathrm{old}}(k)$,
each curve against its own before, both in prevailing moneyness.
Panel~(b) shows the two answers.  At the money they coincide to a few
basis points: \MacSsrRepAtmLogBp\ vol bp for the log-strike field,
\MacSsrRepAtmDlrBp\ for the dollar-strike one, both the half-rule's
$2s_0H$ in action.  (Even the small difference is explainable rather
than numerical: the dollar field's exact ATM answer is
$b(e^{H}-1)$ rather than $bH$, a distinction of second order in the
move.)  Away from the money the two answers have visibly different
\emph{shapes}.  The log-strike field answers with an almost uniform
lift --- the flat blue line: its smile slides up as a block, which is
precisely what the linear $\mathcal{R}=2$ transport
\eqref{eq:ssrshift} builds.  The dollar-strike field answers with a
tilt --- the violet curve, rising from the put side to the call side
--- and the tilt is worth deriving, because a naive read of the
expansion \eqref{eq:ssrellexp} points the wrong way: the displacement
$(1+e^{-k})H$ is \emph{largest} on the put wing, yet the response
there is \emph{smallest}.  Two lines resolve it.  By the midpoint
rule, the dollar-straight field's own smile today is
\[
\sigma_{\mathrm{old}}(k)\approx 0.20+\tfrac{b}{2}\big(e^{k}-1\big),
\qquad\text{with local slope}\qquad
\sigma_{\mathrm{old}}'(k)\approx \tfrac{b}{2}\,e^{k} :
\]
straight in dollars is curved in log-moneyness, and the slope fades
on the put side.  The response at a fixed label is displacement times
slope,
\[
\sigma_{\mathrm{old}}'(k)\,(1+e^{-k})\,H
\;\approx\;\tfrac{b}{2}\,e^{k}\,(1+e^{-k})\,H
\;=\;\tfrac{b}{2}\,(1+e^{k})\,H ,
\]
so the large displacement lands on the flat part of the curve and
the product tilts the \emph{other} way: proportional to $1+e^{k}$,
smallest on the put side --- the violet shape, and exactly what the
relabeling \eqref{eq:ssrell} builds when applied to this smile.  (At
$k=0$ the product is $bH=2s_0H$ again.)  By $k=\pm0.30$ the two
fields' answers are \MacSsrRepSepPutBp\ and \MacSsrRepSepCallBp\ vol
bp apart.  One note of honesty on those two numbers: the first-order
logic just traced puts the separation near \MacSsrRepSepPutBp\ vol bp
on \emph{both} wings; the call side's excess is each field's own
second-order response to a move this large --- the reprice keeps
score exactly where the expansions stop.

So the tie from \cref{sec:ssrhalf} dissolves rather than resolves:
\emph{each} closed form is the exact first-order answer --- for a
different field.  Two fields that today's quotes can barely tell apart
give measurably different tomorrows in the wings.  The conclusion
deserves plain words, because it is the chapter's deepest point.
\textbf{At the money, the frozen-field answer is universal}: twice the
skew, no matter the field's shape --- that is the half-rule, and every
construction in this chapter reproduces it.  \textbf{In the wings, the
answer belongs to the field}, and Chapter~4 showed the quotes do not
pin the field down there (\cref{sec:lvinfo}).  The missing derivative, chased into the
one regime that promised to derive it, reappears one level deeper: the
model answers, but only as precisely as the data determined the model.

\medskip

The frozen field also completes the dial in a satisfying way.  All
three named regimes --- not just $\mathcal{R}=2$ --- can be written as
one rule about \emph{how the field follows the spot}.  Let the field's
strike axis be relabeled, per unit move, by
\begin{equation}
x\;\longmapsto\;x-\tfrac{\mathcal{R}}{2}\,H
\qquad\text{in prevailing log-moneyness }x=\log y .
\label{eq:ssrgrid}
\end{equation}
For $\mathcal{R}=2$ this is $x\mapsto x-H$: exactly the label change
of a field frozen in absolute strike (the quote rule of
\cref{def:ssrsigns}, applied to the field's nodes).  For
$\mathcal{R}=0$ the labels do not move: the field rides with the
forward.  For $\mathcal{R}=1$ it follows halfway.  The consistency
check takes two lines through the half-rule.  Relabeling by
\eqref{eq:ssrgrid} means the field seen in prevailing coordinates is
$\sqrt{v}_{\mathrm{new}}(x)=\sqrt{v}_{\mathrm{old}}
(x+\tfrac{\mathcal{R}}{2}H)$, so the new ATM implied volatility ---
the field's value at the new money, by the chord average --- is
\[
\sigma_{\mathrm{atm}}^{\mathrm{new}}
\approx\sigma_{0}+b\,\tfrac{\mathcal{R}}{2}H
=\sigma_{0}+\mathcal{R}\,s_0\,H ,
\]
using $b=2s_0$.  One relabeling rule for the field generates the
entire dial of \cref{sec:ssrdial}; the stickiness debate is a debate
about how many halves of the move the generator follows.

\medskip

Finally, make the section's title quantitative: measure the ratio as
an \emph{output}.  Define the realized ratio of a reprice,
\[
\widehat{\mathcal{R}}(\tau)
=\frac{\sigma_{\mathrm{atm}}^{\mathrm{new}}(\tau)
-\sigma_{\mathrm{atm}}^{\mathrm{old}}(\tau)}{s_0(\tau)\,H},
\]
the hat marking, as in Chapter~6, a measured estimate: everything on
the right comes out of the pricer.  Panel~(c) of
\cref{fig:ssrreprice} plots it across maturities from $0.05$ to $1.5$
years.  For the straight log-strike field at
$H=-\MacSsrRepRatioMovePct\%$, the curve is a ruler:
$\widehat{\mathcal{R}}$ from \MacSsrRepFlatShort\ at the shortest
maturity to \MacSsrRepFlatLong\ at the longest, with a total spread of
\MacSsrRepFlatSpread\ across the whole range.  For this construction,
$2$ is not merely a short-maturity limit --- at first order in the
skew the frozen field pays out double at \emph{every} maturity, and
the measurement confirms it to the third digit.  The
maturity-independence has a one-sentence reason: at first order the
numerator (the ATM response) is an average of the field's slope over
the neighbourhood the paths visit around the money, the denominator
(the skew) is \emph{half the same average}, and the quotient cancels
whatever maturity dependence the average carries.  Then bend the
field:
add curvature $c\,(\log y)^2$ with $c=0.5$ to the straight generator.
The ratio lifts off $2$ --- to \MacSsrRepBentMean\ at
$H=-\MacSsrRepRatioMovePct\%$ and \MacSsrRepBentBigMean\ at
$H=-4\%$ --- and stays flat in maturity.  The lift is not mysterious;
it is one more chord average.  At the new money the bent generator
reads $\sigma_{0}+bH+cH^{2}$, while the bent field's implied skew is
still $b/2$ --- the curvature term's chord average is $ck^{2}/3$,
whose slope vanishes at the money --- so
\[
\widehat{\mathcal{R}}
=\frac{bH+cH^{2}}{(b/2)\,H}
=2+\frac{2c}{b}\,H ,
\]
and the dotted guides in the panel sit at exactly these levels:
predicted \MacSsrRepBentPred\ and \MacSsrRepBentBigPred, measured
\MacSsrRepBentMean\ and \MacSsrRepBentBigMean.  The ratio is a genuine
output: it reports the field's shape (through $c/b$) and the size of
the move (through $H$), neither of which the constant ``$2$'' of the
regime's name can see.

One honest boundary closes the section.  The frozen field is itself a
hypothesis about the world, and empirically it is not quite right:
measured markets move their smiles by \emph{less} than the frozen
field predicts --- short-dated skew-stickiness ratios sit between the
sticky-strike and sticky-local-vol answers \citep{Bergomi2016}.  The
reprice is the exact consequence of an approximate premise.  That is
why this regime is one setting of the dial and not the dial's
replacement: it answers precisely, and the desk must still decide
whether to believe the premise.
